%_____________________________________________________________________________________________
% Chapter 2: Literature Survey
%_____________________________________________________________________________________________

\chapter{Literature Survey}
\label{ch:lit}

We review concept drift and dataset shift, drift detection, label-free accuracy estimation, delayed labels and pre-registration.

\section{Concept Drift and Dataset Shift}
Gama et al.~\cite{gama2014survey} and Lu et al.~\cite{lu2019review} survey concept drift: changes over time in the joint distribution $P(X, Y)$ of inputs $X$ and labels $Y$. It is useful to separate two cases~\cite{quinonero2009shift}:
\begin{itemize}
\item \textbf{Covariate shift}: $P(X)$ changes while $P(Y\mid X)$ stays the same. The model's decision rule is still correct where it was correct, but the data may move into regions where the model was weak.
\item \textbf{Concept shift}: $P(Y\mid X)$ changes, so the same input now has a different correct label. No input-only signal can see this; only labels reveal it.
\end{itemize}
Label shift, where $P(Y)$ changes and $P(X\mid Y)$ stays the same, is treated by Lipton et al.~\cite{lipton2018labelshift}.

\section{Drift Detection}
\subsection{Detectors that need labels}
Classical detectors monitor the error stream of a model. DDM~\cite{gama2004ddm} signals a warning and then a drift when the error rate rises significantly above its historical minimum. ADWIN~\cite{bifet2007adwin} keeps an adaptive window and cuts it when two sub-windows have significantly different means. Sequential tests such as CUSUM~\cite{page1954cusum} accumulate evidence of a persistent change. All of these need the label of each prediction, so with delayed labels they react late by construction.

\subsection{Label-free detectors}
Label-free detection compares the distribution of the current inputs with a reference sample. Common statistics are the two-sample Kolmogorov--Smirnov (KS) test per feature~\cite{massey1951ks}, the population stability index (PSI) used in credit scoring, divergences between histograms, and the kernel maximum mean discrepancy (MMD)~\cite{gretton2012mmd}. Rabanser et al.~\cite{rabanser2019failing} compared such detectors systematically and noted an important limitation: a detected shift need not reduce accuracy. Open-source monitoring tools such as Evidently~\cite{evidently} and NannyML~\cite{nannyml_cbpe} package these statistics for practitioners.

\section{Label-Free Performance Estimation}
A second line of work estimates the model's accuracy directly, without labels:
\begin{itemize}
\item \textbf{Confidence-based estimation (CBPE)}~\cite{nannyml_cbpe} uses the model's predicted probabilities as the expected accuracy. It is exact only if the model is calibrated on the new data.
\item \textbf{Average thresholded confidence (ATC)}~\cite{garg2022atc} learns a confidence threshold on source data and predicts the target accuracy as the share of target examples above it.
\item \textbf{Difference of confidences (DoC)}~\cite{guillory2021doc} regresses the accuracy change on the change in average confidence.
\item \textbf{Disagreement}~\cite{jiang2022disagreement}: the rate at which two independently trained models disagree tracks their error rate under certain conditions.
\item \textbf{Importance weighting}~\cite{shimodaira2000,sugiyama2007iwcv}: under covariate shift, reweighting labeled source examples by the density ratio $p_{target}(x)/p_{source}(x)$ gives an unbiased estimate of any target expectation, including accuracy. The density ratio can be estimated with a probabilistic classifier that separates source from target samples.
\end{itemize}
Our harm tier (Section~\ref{sec:harm}) uses importance weighting with a logistic domain classifier, applied to the model's correctness on the labeled reference set.

\section{Delayed Labels}
Learning and evaluation with delayed or partially missing labels are studied by Grzenda et al.~\cite{grzenda2020delayed} and Umer and Polikar~\cite{umer2020evl}. In the monitoring setting the key point is causality: at batch $t$, only the labels of batches $b \le t-\max(L,1)$ may be used. Any evaluation that ignores this leaks future information and overstates how early a monitor can warn. Our replay protocol enforces this rule and tests it (Section~\ref{sec:tests}).

\section{Pre-Registration and Reproducibility}
Pre-registration means fixing hypotheses, analysis methods and success criteria before seeing the data that will judge them~\cite{nosek2018prereg}. It reduces the freedom to choose, after the fact, the analysis that looks best. Van Miltenburg et al.~\cite{vanmiltenburg2021prereg} discuss pre-registration for machine learning and natural language processing research. Sculley et al.~\cite{sculley2015debt} and Breck et al.~\cite{breck2017mltest} argue that production ML systems need explicit, testable monitoring. We apply pre-registration to the evaluation of a monitoring system, including its alert thresholds, and we use public datasets from OpenML~\cite{vanschoren2014openml} so that every result can be reproduced.

\section{Summary and Research Gap}
\tabref{tab:lit} summarises the main families of methods.

\begin{table}[H]
\caption{Families of drift-monitoring methods and their limits in the delayed-label setting.}
\label{tab:lit}
\centering
\small
\begin{tabularx}{\textwidth}{@{}p{3.4cm}p{3.2cm}X@{}}
\toprule
Family & Examples & Limitation for early warning with delayed labels \\
\midrule
Error-stream detectors & DDM, ADWIN, CUSUM & Need labels; with delay $L$ they react at least $L$ batches late \\
Input-distribution tests & KS, PSI, MMD, fused scores & Detect any change, harmful or not; thresholds in dataset-specific units \\
Confidence-based estimators & CBPE, ATC, DoC & Assume calibration on new data; can stay confident while wrong \\
Disagreement & auxiliary models & Moves with any shift that the auxiliary models notice \\
Importance weighting & density-ratio reweighting & Valid only under covariate shift; needs a labeled reference sample \\
\bottomrule
\end{tabularx}
\end{table}

Three things are missing from this work:
\begin{enumerate}
\item Most evaluations report whether drift is detected, not whether a warning comes \emph{before} the accuracy drop or \emph{before} delayed labels could reveal it.
\item Alert thresholds are rarely tested on datasets other than the ones they were chosen on.
\item Few studies measure, on real data with known ground truth, how often a drift alarm corresponds to actual harm.
\end{enumerate}
The project addresses these gaps with a benchmark in which harm is measured, warning tiers defined relative to label arrival, and pre-registered confirmation on datasets never used during design.
